\documentclass[11pt,a4paper]{article}
\usepackage[T2A]{fontenc}
\usepackage[utf8]{inputenc}
\usepackage[russian]{babel}
\usepackage{amsmath,amssymb}
\usepackage[margin=20mm,headheight=15pt]{geometry}
\usepackage{xcolor,enumitem,needspace,fancyhdr,hyperref}
\definecolor{accent}{HTML}{174C72}
\definecolor{muted}{HTML}{59636C}
\hypersetup{colorlinks=true,linkcolor=accent,urlcolor=accent,
 pdftitle={Архив экзаменационных задач по дифференциальным уравнениям}}
\setlength{\parindent}{0pt}
\setlength{\parskip}{5pt}
\setlist[itemize]{nosep,leftmargin=1.5em}
\pagestyle{fancy}
\fancyhf{}
\fancyhead[L]{\small\color{muted}Дифференциальные уравнения}
\fancyhead[R]{\small\color{muted}Архив заданий}
\fancyfoot[C]{\small\thepage}
\renewcommand{\headrulewidth}{0.3pt}
\newcommand{\sect}[1]{\Needspace{85mm}\section*{\color{accent}#1}}
\newcommand{\task}[2]{\par\vspace{6pt}
 {\bfseries\color{accent}#1. #2}\par}
\newcommand{\note}[1]{\par{\small\color{muted}#1}\par}
\newcommand{\R}{\mathbb R}
\begin{document}

{\huge\bfseries\color{accent}Архив экзаменационных задач}\par
{\large По дифференциальным уравнениям}\par
\vspace{4pt}

\sect{1. Линейные уравнения: свойства и существование}

\begin{minipage}[t]{\linewidth}
\task{1}{Ограниченность всех решений}
При каких значениях параметров $a,b$ все решения уравнения
\[
 y''+ay'+by=0
\]
ограничены?
\end{minipage}\par\vspace{6pt}

\begin{minipage}[t]{\linewidth}
\task{2}{ЛОУ второго порядка с заданными решениями}
Существует ли линейное однородное уравнение второго порядка с вещественными
постоянными коэффициентами, имеющее частные решения $\cos x$ и $2$?
Ответ обоснуйте.
\end{minipage}\par\vspace{6pt}

\begin{minipage}[t]{\linewidth}
\task{3}{Минимальный порядок ЛОУ}
Каков минимальный порядок линейного однородного уравнения с вещественными
постоянными коэффициентами, имеющего частные решения $\cos x$ и $2$?
Ответ обоснуйте.
\end{minipage}\par\vspace{6pt}

\begin{minipage}[t]{\linewidth}
\task{4}{ЛНУ с заданными решениями}
Существует ли линейное неоднородное уравнение с вещественными постоянными
коэффициентами, имеющее частные решения $\cos x$ и $2$? Ответ обоснуйте.
\end{minipage}\par\vspace{6pt}

\sect{2. Уравнения и задачи Коши}

\begin{minipage}[t]{\linewidth}
\task{6}{ЛНУ с постоянными коэффициентами}
Решить уравнение
\[
 y''+3y'+2y=\frac{1}{e^x+1}.
\]
\end{minipage}\par\vspace{6pt}

\begin{minipage}[t]{\linewidth}
\task{7}{Неоднородное уравнение с тригонометрией}
Решить уравнение
\[
 y''+y=\frac{1}{\sin x}.
\]
\end{minipage}\par\vspace{6pt}

\begin{minipage}[t]{\linewidth}
\task{8}{Уравнение без независимой переменной}
Решить уравнение
\[
 y''=\frac{1}{y^3}.
\]
\end{minipage}\par\vspace{6pt}

\begin{minipage}[t]{\linewidth}
\task{9}{Задача Коши с котангенсом}
\[
 \begin{cases}
 y''+(y')^2\operatorname{ctg}y=0,\\
 y(0)=\dfrac{\pi}{2},\qquad y'(0)=1.
 \end{cases}
\]
\end{minipage}\par\vspace{6pt}

\begin{minipage}[t]{\linewidth}
\task{11}{Задача Коши из текстовой реплики}
\[
 \begin{cases}
 y^2y''+(y'-1)^2=1,\\
 y(0)=1,\qquad y'(0)=2.
 \end{cases}
\]
\end{minipage}\par\vspace{6pt}

\begin{minipage}[t]{\linewidth}
\task{12}{Единственность по теореме Коши}
\[
 y^2y''+y'(y'-1)^2=0.
\]
По теореме Коши доказать единственность и найти решение.
\note{Начальные условия не сохранились. Ответ: $y=x+1$.}
\end{minipage}\par\vspace{6pt}

\sect{3. Системы и положения равновесия}

\begin{minipage}[t]{\linewidth}
\task{13}{Линейная система из контрольной}
\[
 \begin{cases}
 \dot x=-2x-y+z,\\
 \dot y=5x-y+4z,\\
 \dot z=5x+y+2z.
 \end{cases}
\]
Решить систему. Исследовать нулевое решение на устойчивость.
Есть ли неустойчивые решения системы? Указать хотя бы одно, если да.
\end{minipage}\par\vspace{6pt}

\begin{minipage}[t]{\linewidth}
\task{14}{Равновесия экспоненциальной системы}
Являются ли точки $(0,0)$ и $(0,1)$ положениями равновесия системы
\[
 \begin{cases}
 \dot x=e^{x+2y}-\cos 3x,\\
 \dot y=\sqrt{4+8x}-2e^y.
 \end{cases}
\]
Для каждой точки, являющейся положением равновесия, определить тип и
исследовать устойчивость.
\end{minipage}\par\vspace{6pt}

\begin{minipage}[t]{\linewidth}
\task{15}{Равновесия системы с синусом}
Являются ли точки $(0,0)$, $(\pi,0)$ и $(\pi,1)$ положениями равновесия системы
\[
 \begin{cases}
 \dot x=y,\\
 \dot y=\sin(x+y).
 \end{cases}
\]
Для каждого положения равновесия определить тип и исследовать устойчивость.
\end{minipage}\par\vspace{6pt}

\begin{minipage}[t]{\linewidth}
\task{16}{Система из рукописного списка}
\[
 \begin{cases}
 \dot x=8\sin^2 y,\\
 \dot y=(x-2)(x^2+x).
 \end{cases}
\]
\note{На листе отмечены точки $(2,0)$ и $(0,0)$.
Формулировка вопроса не сохранилась.}
\end{minipage}\par\vspace{6pt}

\sect{4. Устойчивость и качественное исследование}

\begin{minipage}[t]{\linewidth}
\task{17}{Устойчивость решения задачи Коши}
Исследовать на устойчивость по Ляпунову решение задачи Коши
\[
 \begin{cases}
 (t+1)x'+tx=0,\\
 x(0)=0.
 \end{cases}
\]
\end{minipage}\par\vspace{6pt}

\begin{minipage}[t]{\linewidth}
\task{18}{Равновесия и фазовый портрет (№ 888)}
Найти положения равновесия системы
\[
 \begin{cases}
 \dot x=-y\cos x,\\
 \dot y=\sin x.
 \end{cases}
\]
Пользуясь теоремой Пуанкаре--Перрона, установить, устойчивы ли они.
Построить фазовый портрет.
\end{minipage}\par\vspace{6pt}

\begin{minipage}[t]{\linewidth}
\task{19}{Число решений через прямую}
Для уравнения
\[
 y'=\frac{e^{x+1}}{x-1}
\]
найти, сколько решений проходит через прямую $x=1$.
\end{minipage}\par\vspace{6pt}

\begin{minipage}[t]{\linewidth}
\task{20}{Вертикальная асимптота}
Для уравнения
\[
 y'=e^y(y-1)
\]
выяснить, имеет ли интегральная кривая вертикальную асимптоту $x=C$
при $y\to+\infty$.
\end{minipage}\par\vspace{6pt}

\begin{minipage}[t]{\linewidth}
\task{21}{Параметр и седло}
При каких значениях параметра $a$ получается седло?
\[
 \begin{cases}
 \dot x=ax+ay,\\
 \dot y=a^2y^2.
 \end{cases}
\]
\end{minipage}\par\vspace{6pt}

\begin{minipage}[t]{\linewidth}
\task{25}{Устойчивость по определению (№ 881)}
Пользуясь определением устойчивости по Ляпунову, выяснить, устойчивы ли
решения данных уравнений с указанными начальными условиями:
\[
\begin{aligned}
 \text{а)}\quad &3(t-1)\dot x=x,       &&x(2)=0;\\[5pt]
 \text{б)}\quad &\dot x=4x-t^2x,      &&x(0)=0;\\[5pt]
 \text{в)}\quad &\dot x=t-x,          &&x(0)=1;\\[5pt]
 \text{г)}\quad &2t\dot x=x-x^3,      &&x(1)=0.
\end{aligned}
\]
\end{minipage}\par\vspace{6pt}

\sect{5. Неполные условия и упоминания типов задач}

\begin{minipage}[t]{\linewidth}
\task{26}{Частное решение второго вида}
Существует ли частное решение второго вида ЛНУ второго порядка?
\end{minipage}\par\vspace{6pt}

\vspace{5pt}
\textbf{Ещё упомянуты следующие задания без полных условий:}
\begin{itemize}
 \item Уравнение Эйлера.
 \item Первая задача из последней контрольной --- решить систему
       (была почти у всех).
 \item Уравнение с отношением $y/x$; по воспоминанию, интегрирование
       требовало двух замен, одна из которых связана с тригонометрией.
 \item Вторая задача на интегральную формулу Коши.
 \item Линейное неоднородное уравнение с постоянными коэффициентами.
\end{itemize}

\end{document}
